\documentclass{beamer}
\usepackage{amsfonts,amsmath,oldgerm}

\usepackage[utf8]{inputenc}
\usepackage[T1]{fontenc}
\usepackage{graphicx}
\usepackage{listings}
\usepackage{hyperref}
\usepackage{verbatim}


\usepackage[export]{adjustbox}

\usetheme{sintef}

\newcommand{\testcolor}[1]{\colorbox{#1}{\textcolor{#1}{test}}~\texttt{#1}}

\usefonttheme[onlymath]{serif}

\titlebackground*{assets/background}

\newcommand{\hrefcol}[2]{\textcolor{cyan}{\href{#1}{#2}}}

\title{Big data - Parte 2}

\author{\href{mailto:oscar.bustos@javeriana.edu.co}{Oscar Bustos}}

\begin{document}
\maketitle


\begin{frame}

This template is a based on \hrefcol{https://www.overleaf.com/latex/templates/sintef-presentation/jhbhdffczpnx}{SINTEF Presentation} from \hrefcol{mailto:federico.zenith@sintef.no}{Federico Zenith} and its derivation \hrefcol{https://github.com/TOB-KNPOB/Beamer-LaTeX-Themes}{Beamer-LaTeX-Themes} from Liu Qilong and Andrea Gasparini

% This template is released under \hrefcol{https://creativecommons.org/licenses/by-nc/4.0/legalcode}{Creative Commons CC BY 4.0} license

\end{frame}

\section{Introducción}

\begin{frame}{¿Qué es el Procesamiento de Lenguaje Natural?}
  \begin{itemize}
    \item El Procesamiento de Lenguaje Natural (PLN) es una subdisciplina de la inteligencia artificial que se enfoca en la interacción entre las computadoras y los lenguajes humanos.
    \item Implica el desarrollo de algoritmos y modelos para comprender, interpretar y generar lenguaje humano.
  \end{itemize}
\end{frame}

\section{Índice Invertido}

\begin{frame}{Índice Invertido}
  \begin{itemize}
    \item Un índice invertido es una estructura de datos utilizada en sistemas de recuperación de información.
    \item Mapea términos (palabras) a sus ubicaciones en un conjunto de documentos o textos.
    \item Esto permite búsquedas de texto completo eficientes y es un componente fundamental de los motores de búsqueda.
  \end{itemize}
\end{frame}


\section{Embeddings de Palabras}

\begin{frame}{Embeddings de Palabras}
  \begin{itemize}
    \item Los embeddings de palabras son representaciones vectoriales de palabras en un espacio vectorial continuo.
    \item Capturan relaciones semánticas entre palabras y se utilizan en diversas tareas de PLN como el análisis de sentimientos, la traducción automática y más.
  \end{itemize}
\end{frame}

\begin{frame}[fragile]{Verbatim Code Example}
    \begin{verbatim}
    #include <stdio.h>

    int main() {
        printf("Hello, world!\n");
        return 0;
    }
    \end{verbatim}
\end{frame}



\section{Conclusiones}

\begin{frame}{Conclusiones}
  \begin{itemize}
    \item El Procesamiento de Lenguaje Natural es un campo diverso y emocionante con muchas aplicaciones.
    \item Los índices invertidos permiten una recuperación de texto eficiente.
    \item Los embeddings de palabras capturan el significado semántico y son cruciales para varias tareas de PLN.
  \end{itemize}
\end{frame}

\begin{frame}{¿Preguntas?}
  \begin{center}
    ¡Gracias por su atención!\\
    ¿Alguna pregunta?
  \end{center}
\end{frame}




\end{document}
